\section{Introduction}


% ¶1: Context — agents are powerful and widely deployed
AI coding agents have become widely deployed tools for software
development~\cite{agentsight}. They operate as long-running processes with
access to shells, source trees, package managers, and external APIs,
autonomously writing code, running tests, and managing infrastructure across
extended sessions.

% ¶2: Agents need deterministic harnesses to apply behavioral policies
As agents gain autonomy, they require harnesses to apply behavioral
policies. Projects encode such policies in instruction files such
as \texttt{CLAUDE.md} and
\texttt{AGENTS.md}~\cite{chatlatanagulchai2025manifests,lulla2026agentsmd}.
Our empirical study of 64~projects extracts 2{,}116 statements and finds
that 64\% are behavioral directives, 83\% of which involve system-level
behavior, and 81\% of projects require cross-event state tracking
(\S\ref{sec:empirical}).
A directive is \emph{per-event} if it can be checked against a single
operation (``do not run \texttt{git push --force}'') and
\emph{cross-event} if it requires state across operations (``run tests
before committing''---the commit is forbidden unless a test execution
has been observed since the last source change).
Because LLM compliance is inherently
probabilistic~\cite{liu2024lost}, policy cannot reside inside the
model~\cite{langchain-harness,fowler-harness}: even cooperative agents
routinely violate declared constraints through long-horizon planning
errors, shell-script side effects, and opaque tool
behavior~\cite{jimenez2024swebench}.

% ¶3: Gap — context gap + semantic gap
Enforcing these policies faces two compounding gaps. First, a
\emph{context gap}: 74\% of system-level directives are not
self-contained---they reference project domains (``run the test suite'')
or delegate judgment to the current task (``unless explicitly
requested'') that cannot be resolved without the agent's own project
and session context (\S\ref{sec:empirical}). A static policy authored
by an external administrator cannot supply this context.

Second, a \emph{semantic gap} between declared policy and system effects.
Prompt constraints remain probabilistic, while tool-level
guards~\cite{wang2025agentspec,shi2025progent} lose coverage when agents
shell out or spawn subprocesses: the same \texttt{git~push} can be reached
through a tool call, a \texttt{bash~-c} string, or a compiled binary.
OS-level policy~\cite{ciliumtetragon,flume,pasquier2018camquery} sees
these effects but requires policies to be statically pre-written and
returns opaque errors with no connection to declared intent.
The agent---the entity that holds both the project context and the
task intent needed to instantiate 74\% of its own rules---has no
programmatic interface to define or evolve the constraints governing
its system-level actions.

% ¶4: Insight — programmable interface, agent as control plane
Our insight is that the policy engine should be exposed as a
programmable interface that lets agents become the control plane of
their own system-level actions. Because the agent is the only entity
that holds the project structure and current task intent needed to
concretize its rules, the agent must author its own enforcement
policy---trusted to honestly declare intent, but not trusted to
reliably follow it during execution.

% ¶5: System — ActPlane
ActPlane is such an interface. Agents and developers declare behavioral
policies in a compact DSL that agents themselves can generate from
natural-language directives; ActPlane compiles them into labeled
information-flow rules and loads them into an eBPF/BPF-LSM backend. Named
labels propagate across process, file, and network objects at kernel hooks,
encoding execution history as checkable state: when a process reads a file
labeled \texttt{SECRET}, the process acquires that label; when it forks, the
child inherits it. When a labeled subject matches a rule, ActPlane returns
semantic feedback---which rule was matched, why, and what alternative path
satisfies it. Each rule independently selects notify (observe), block
(pre-operation denial), or kill (post-operation termination).
Policies are layered: higher-authority rules (e.g., platform-level
``do not leak secrets'') cannot be weakened by the agent, while
agent-authored rules can be revised as task intent evolves.

This paper makes three contributions.

\begin{enumerate}
\item An \textbf{empirical study} of instruction files from 64~popular
projects, showing that cross-event information-flow policies are pervasive
(81\% of repositories) and that 74\% of system-level directives require
project or task context that only the agent can supply.

\item \textbf{ActPlane}, a programmable OS-level control plane for agent
harnesses. ActPlane compiles behavioral policies into eBPF/BPF-LSM labeled
information-flow rules, observes and enforces them across process, file,
and network boundaries, returns semantic feedback that reconnects rule
matches to intent-level remediation, and supports layered policies where
higher-authority constraints cannot be weakened by the agent.

\item An \textbf{evaluation} of ActPlane's policy coverage, runtime
cost, and feedback effect. We measure coverage over 607~OS-level
directives from the empirical study, quantify per-event and end-to-end
overhead, and evaluate repository-grounded checklist satisfaction on a
20-task OS-observable OctoBench subset against prompt-only and tool-regex
baselines.
\end{enumerate}
